\chapter{Worked example: commitment schemes}

Following Rosulek, \emph{The Joy of Cryptography} (commitments). A commitment binds
a committer to a value while hiding it until opening. Perfect hiding and perfect
binding are information-theoretically incompatible, so the two properties are
shown on complementary schemes: a mask (one-time-pad) commitment is perfectly
hiding, and the identity commitment is perfectly binding.

\begin{definition}[Commitment scheme]
  \label{def:commscheme}
  \uses{def:spcomp}
  \lean{CatCrypt.Examples.Commitment.CommScheme}
  \leanok
  A commitment scheme has a $\mathsf{commit}$ map (value and randomness to a
  commitment) and a $\mathsf{verify}$ check used at opening.
\end{definition}

\begin{theorem}[Perfect hiding of the mask commitment]
  \label{thm:comm-hiding}
  \uses{def:commscheme, thm:otp-perfect}
  \lean{CatCrypt.Examples.Commitment.maskComm_perfect_hiding}
  \leanok
  The mask commitment $\mathsf{commit}(r, m) = r \oplus m$ with uniform randomness
  $r$ is perfectly hiding: the hiding advantage is exactly $0$ for every
  adversary, by the one-time-pad argument.
\end{theorem}

\begin{theorem}[Perfect binding of the identity commitment]
  \label{thm:comm-binding}
  \uses{def:commscheme}
  \lean{CatCrypt.Examples.Commitment.idComm_perfectly_binding}
  \leanok
  The identity commitment cannot be opened to two distinct messages --- it is
  perfectly binding (and, being the complement of the mask commitment, not
  hiding).
\end{theorem}
